\label{exp:cem02}

\subsubsection*{Постановка}
\textbf{Физический вопрос.}
Постоянная тонкой структуры $\alpha$ и масштаб $v$ в Стандартной модели обычно \emph{входят} как параметры.
В модели coupling и массы upstairs должны вытекать из геометрии насыщающей фазы и лестницы §7--§8 без кругового «подставили $\alpha$ из CODATA, получили $\alpha$''.

\textbf{Что проверяем.}
\begin{enumerate}[label=\textbf{(\alph*)},leftmargin=*,itemsep=0.25em]
\item $\alpha^{-1}_{\mathrm{pref}}$ из геометрического счёта ($M=97$ и связанные факторы, \cref{ch:alpha}) согласуется с CODATA как \emph{выход}, не как fit.
\item Массы $m_H$, $m_p$, $m_e$ (и связанные якоря) из upstairs-лестницы дают относительные отклонения $|\delta m/m|$ к PDG в заявленной полосе модели.
\item Независимые косвенные наблюдаемые не раздувают $E_n$ сверх порога сравнения (\cref{sec:compare-reference}).
\end{enumerate}

\textbf{Рабочая гипотеза.}
При фиксированной лестнице и SI-мосте CE-M-00 относительные ошибки масс остаются на уровне $10^{-3}$--$10^{-4}$, а $\alpha^{-1}$ попадает в полосу CODATA без дополнительных свободных параметров.

\textbf{Наблюдаемые.}
$\alpha^{-1}_{\mathrm{pref}}$, $v$, $|\delta m_H/m_H|$, $|\delta m_p|$, $|\delta m_e|$ (табл.~\ref{tab:exp-ce-m-02}).

\textbf{Критерий согласия с теорией.}
$|\delta m/m|$ в полосе, зафиксированной для upstream-якорей; $\alpha^{-1}$ согласован с экспериментом как следствие геометрии; $E_n\lesssim 1$ для независимых косвенных.

\subsubsection*{Теоретическая часть}
Глава~\cref{ch:alpha} связывает $\alpha_{\mathrm{pref}}$ с фазой насыщения и дискретным счётом степеней свободы; глава~\cref{ch:si-sm} задаёт массы через лестницу и $v$.
Массы на~$T$ --- \emph{косвенные}: они зависят от выбранных upstairs-якорей (\cref{sec:multiple-observables}); сравнение с PDG всегда сопровождается явным перечнем зависимостей.

\subsubsection*{Ход эксперимента}
Воспроизведение: \texttt{python verify\_principles.py --suite alpha --device cpu}.

\begin{enumerate}[label=\textbf{Шаг \arabic*:},leftmargin=*,itemsep=0.35em]
\item Зафиксировать коммит; прогнать suite \texttt{alpha}.
\item Снять $\alpha^{-1}_{\mathrm{pref}}$ и $v$ из SI-цепочки.
\item Вычислить $|\delta m/m|$ для $H$, $p$, $e$ относительно PDG при той же лестнице.
\end{enumerate}

\begin{table}[htbp]
\centering
\small
\caption{SI-мост и массы upstairs (CE-M-02)}
\label{tab:exp-ce-m-02}
\begin{tabular}{@{}ll@{}}
\toprule
Поле & Значение \\
\midrule
$\alpha^{-1}_{\mathrm{pref}}$ & $137{,}03599919$ \\
$v$ [GeV] & $246{,}0873$ \\
$|\delta m_H/m_H|$ & $4{,}73\times 10^{-4}$ \\
$|\delta m_p|$ & $3{,}16\times 10^{-3}$ \\
$|\delta m_e|$ & $4{,}50\times 10^{-3}$ \\
\bottomrule
\end{tabular}
\end{table}

\subsubsection*{Анализ, расчёт погрешности и выводы}
$\alpha^{-1}_{\mathrm{pref}}$ совпадает с CODATA в пределах заявленной точности модели --- coupling не подставлен задним числом.
Ошибки масс --- относительные $|\delta|$ к PDG при \emph{одной} фиксированной лестнице; это не независимые «пять разных экспериментов'', а согласованный блок (\cref{sec:multiple-observables}).

Погрешность: тип~B для алгебры $\alpha$; массы --- сравнение с внешним эталоном PDG (тип~B) плюс учёт зависимых якорей.

\textbf{Вывод.}
Цепочка §7--§8.2 замкнута: геометрия $\to$ $\alpha$, лестница $\to$ $v$ и массы без fit $\alpha$ из эксперимента.
\textbf{Статус записи:} опыт закрыт по табл.~\ref{tab:exp-ce-m-02} и порогам \cref{sec:compare-reference}.
